
\chapter{Robot-Side Integration}

\section{FastAPI Application and Service Composition}
\begin{itemize}
    \item Describe the server startup process, application state, and service initialization.
    \item Explain how configuration loading leads to either an in-process perception pipeline or worker-backed runtime.
    \item Present the main route groups: detection, captioning, chat, vision chat, memory, config, worker, and dashboard.
\end{itemize}

\section{Detection Subsystem}
\begin{itemize}
    \item Describe the detection service abstraction and the supported backends: YOLO, RT-DETR, RF-DETR, and OWL-V2.
    \item Explain ontology handling for open-vocabulary detection.
    \item Discuss thresholding, device selection, and detector interchangeability.
\end{itemize}

\section{Persistent Tracking and Identity Maintenance}
\begin{itemize}
    \item Describe the role of \texttt{SceneMemory} in maintaining object continuity across frames.
    \item Explain appearance embedding extraction, crop preparation, and use of a visual encoder for re-identification.
    \item Explain the Hungarian matching procedure and the weighted combination of appearance and geometry.
    \item Cover unmatched-track aging, dormant-track removal, and object hit statistics.
\end{itemize}

\section{Pepper-Specific Fusion with Robot-Native People Perception}
\begin{itemize}
    \item Explain why server-side visual detections alone are insufficient for socially grounded interaction.
    \item Describe geometric projection from robot person angles into image coordinates.
    \item Explain the matching between Pepper people and server person detections using angular similarity, previous bindings, tracker stability, and distance-size consistency.
    \item Describe the generation of synthetic person detections when Pepper reports a person but no visual bounding box matches.
    \item Explain how social attributes such as waving, sitting, looking at the robot, engagement zone, age, gender, smile, and expressions are merged into object state.
\end{itemize}

\section{Set-of-Mark Rendering}
\begin{itemize}
    \item Describe why Set-of-Mark visualization is used for VLM grounding.
    \item Explain bounding boxes, labels, optional masks, optional polygons, and style scaling.
    \item Mention the available mask backends and the role of the overlay in the scene graph prompt pipeline.
\end{itemize}

\section{Scene Graph Generation}
\begin{itemize}
    \item Present the scene graph service as a strategy layer with four modes: rules, VLM, RelTR, and hybrid.
    \item Explain the rule-based backend:
    \begin{itemize}
        \item configurable predicates;
        \item spatial, directional, overlap, containment, and label-pair rules;
        \item automatic color-attribute extraction.
    \end{itemize}
    \item Explain the VLM backend:
    \begin{itemize}
        \item prompt construction from predicates and captions;
        \item structured output parsing and repair;
        \item filtering of hallucinated relations against current detections.
    \end{itemize}
    \item Explain the optional RelTR backend and why it is treated as a complementary relation source.
    \item Explain the hybrid mode as a union of complementary evidence sources.
\end{itemize}

\section{Scene Memory as a Dynamic World Model}
\begin{itemize}
    \item Describe the internal state representation: tracked objects, relationships, captions, Pepper-person bindings, and timestamps.
    \item Explain object fields that matter for dialogue: label, attributes, status, hits, frame IDs, scan IDs, estimated bearing, Pepper person ID, and robot distance.
    \item Explain relationship accumulation, unary-attribute handling, and caption insertion.
    \item Cover memory pruning, TTLs, and configurable object, relation, and caption limits.
\end{itemize}

\section{Captioning and Auxiliary Vision Services}
\begin{itemize}
    \item Describe the captioning subsystem and its role in fast single-frame scene summarization.
    \item Explain that the caption endpoint can launch detection asynchronously so that later dialogue is grounded in updated memory.
    \item Describe the separate \texttt{vision\_chat} endpoint for direct image-plus-question multimodal interaction.
\end{itemize}

\section{Worker Runtime and Runtime Switching}
\begin{itemize}
    \item Describe the worker manager, internal worker API, and lazy startup behavior.
    \item Explain why the worker mode exists: isolate heavy inference, manage warmup, and keep the web process lighter.
    \item Cover health checks, restart logic, circuit breaker behavior, and hot runtime updates.
\end{itemize}

\section{Dashboard and Runtime Administration}
\begin{itemize}
    \item Describe the dashboard pages and their practical roles.
    \item Explain live WebSocket updates for detections, captions, memory, and chat messages.
    \item Cover operator-editable settings: detector backend, thresholds, scene graph mode, VLM prompts, worker controls, memory limits, chat/caption providers, and pipeline presets.
\end{itemize}
